% ============================================================================
\chapter{VMM：虚拟内存与分页}
\label{ch:vmm}
% Stage 6 — 对应 2026-03-11 changelog
% ============================================================================

\section{本章目标}

\begin{itemize}
  \item 理解 x86 两级页表（Page Directory + Page Table）
  \item 实现 \texttt{vmm\_map\_page} / \texttt{vmm\_unmap\_page}
  \item 用 4KiB 页表替换 boot.asm 的 PSE 临时映射，开启正式分页
  \item 实现按需分配页表（\texttt{vmm\_ensure\_page\_table}）
  \item Shell 命令：\texttt{vmm state}、\texttt{vmm lookup}、\texttt{vmm fault}
\end{itemize}

\section{背景知识：x86 二级分页}

虚拟地址（32 位）被 CPU 拆成三段：

\begin{lstlisting}[style=shellstyle]
  31        22 21       12 11        0
 +-----------+-----------+-----------+
 | PD index  | PT index  |  Offset   |
 |  (10 bit) |  (10 bit) |  (12 bit) |
 +-----------+-----------+-----------+
\end{lstlisting}

翻译流程：\reg{CR3} → \texttt{PD[pd\_idx]} → \texttt{PT[pt\_idx]} → 物理页帧 + Offset

\subsection{PDE / PTE 格式}

\begin{bitfield}
  \bit{0}{P (Present)：1=有效，0=访问触发 \#PF}
  \bit{1}{R/W：1=可读写，0=只读}
  \bit{2}{U/S：1=用户态可访问，0=仅内核}
  \bit{5}{A (Accessed)：CPU 访问后自动置 1}
  \bit{6}{D (Dirty)：CPU 写入后自动置 1（仅 PTE）}
  \bit{12--31}{物理页帧地址高 20 位}
\end{bitfield}

\begin{thinkbox}
PDE 和 PTE 都有 R/W 位。如果 PDE 设了 Read-Only 但 PTE 设了 Writable，
最终这个页是可写还是只读？提示：权限取\textbf{交集}（最严格的生效）。
\end{thinkbox}

\section{核心实现}

\subsection{\texttt{vmm\_ensure\_page\_table}：按需分配页表}

\begin{lstlisting}[style=cstyle]
static page_table_t* vmm_ensure_page_table(uint32_t pd_idx) {
    uint32_t pde = g_kernel_pd.entries[pd_idx];
    if (pde & PDE_PRESENT) {
        return (page_table_t*)PHYS_TO_VIRT(pde & ~0xFFFu);
    }
    // 从 PMM 分配一页作为新页表
    uint32_t pt_phys = (uint32_t)pmm_alloc_page();
    page_table_t* pt = (page_table_t*)PHYS_TO_VIRT(pt_phys);
    memzero_page(pt);  // 所有 PTE = 0（未映射）
    g_kernel_pd.entries[pd_idx] = pt_phys | PDE_PRESENT | PDE_WRITABLE;
    return pt;
}
\end{lstlisting}

\begin{whybox}
为什么 PDE 中存物理地址而非虚拟地址？

PDE/PTE 是给 CPU 的 MMU 硬件读的。MMU 工作在物理地址层面——
它拿到 \reg{CR3} 的物理地址，读 PDE 的物理地址找到页表，再读 PTE 的物理地址找到页帧。
我们在 C 代码中访问页表时必须用 \texttt{PHYS\_TO\_VIRT} 转为虚拟地址。
\end{whybox}

\subsection{\texttt{vmm\_map\_page}}

\begin{lstlisting}[style=cstyle]
int vmm_map_page(uint32_t virt, uint32_t phys, uint32_t flags) {
    uint32_t pdi = virt >> 22;
    uint32_t pti = (virt >> 12) & 0x3FF;
    page_table_t* pt = vmm_ensure_page_table(pdi);
    pt->entries[pti] = (phys & ~0xFFF) | (flags & 0xFFF);
    if (g_vmm_ready) vmm_invlpg(virt);  // 刷新 TLB
    return 0;
}
\end{lstlisting}

\begin{thinkbox}
\texttt{vmm\_invlpg} 刷新的是单个虚拟地址在 TLB 中的缓存。
如果我们映射了 1000 个页，需要调用 1000 次 \texttt{invlpg} 吗？
有没有更高效的方法？（提示：重新加载 \reg{CR3} 会刷新整个 TLB。）
\end{thinkbox}

\subsection{\texttt{vmm\_init}：替换 boot PSE 映射}

\begin{enumerate}
  \item 清零页目录
  \item 查询 PMM 确定需要映射的物理范围
  \item 逐页建立 identity mapping + high-half mapping
  \item 切换 \reg{CR3} 到新页目录
  \item 设置 \texttt{g\_next\_vaddr}（bump allocator 水位线）
\end{enumerate}

\section{Bump 分配器：\texttt{vmm\_alloc\_pages}}

\begin{lstlisting}[style=cstyle]
void* vmm_alloc_pages(unsigned count) {
    uint32_t start = g_next_vaddr;
    g_next_vaddr += count * VMM_PAGE_SIZE;
    for (unsigned i = 0; i < count; i++) {
        uint32_t phys = (uint32_t)pmm_alloc_page();
        vmm_map_page(start + i * 4096, phys, PTE_PRESENT | PTE_WRITABLE);
    }
    return (void*)start;
}
\end{lstlisting}

\begin{designbox}
Bump allocator 不能回收虚拟地址空间——\texttt{free\_pages} 只归还物理页。
这是有意为之的简化。Stage 9 的 VMA 红黑树会取代它做正式的虚拟地址管理。
\end{designbox}

\section{验证}

\begin{lstlisting}[style=shellstyle]
szy-kernel > vmm state
VMM: paging enabled
VMM: 512 / 1024 PDE entries present
szy-kernel > vmm lookup 0xC0100000
virt 0xC0100000 -> phys 0x00100000  [P RW S]
szy-kernel > vmm fault
# 触发 Page Fault → 打印 CR2 + error code → halt
\end{lstlisting}

\section{思考与练习}

\begin{exercise}
用 \texttt{vmm pd 768 8} 查看高半区的页目录条目。每个 PDE 覆盖多少虚拟内存？
计算前 8 个 PDE 一共覆盖了多少 MiB。
\end{exercise}

\begin{exercise}
手动用 \texttt{vmm map 0xF0000000 0x00200000} 建立一个映射，然后
\texttt{vmm lookup 0xF0000000} 验证。再 \texttt{vmm unmap 0xF0000000} 后
再查一次。思考：映射后能否直接读写 \hex{F0000000}？（当前 Shell 没有
直接读内存的命令，但你可以在 C 代码中尝试。）
\end{exercise}

\begin{exercise}
TLB 是"页表翻译缓存"。如果修改了 PTE 但没有 \texttt{invlpg}，CPU 可能
使用旧的缓存结果。在 \texttt{vmm\_map\_page} 中注释掉 \texttt{invlpg} 调用，
然后 map → unmap → 重新 map 同一个地址，观察是否出现异常行为。
\end{exercise}

\section{本章小结}

分页开启后，内核的每一次内存访问都受 MMU 控制。
我们有了 4KiB 粒度的映射能力和 Page Fault 诊断。
下一章把内核搬到高地址（High-Half Kernel）。
